% Page 011 — page imprimée 7
% Suite de l'interrogatoire de Jacques Crégut
% Reconstruction éditable en LaTeX.


faisaient corps de garde. Castanet et sa troupe ayant mis le feu au lieu et tué cinq personnes hommes ou femmes et les habitants auraient tué cinq hommes de la troupe de Castanet.

Ils furent au \textbf{Château Daires} dans lequel campait Laroze, lieutenant de Castanet qui fut, dans la nuit avec un détachement de 20 hommes. L’ayant reconnu il le vit avec un justaucorps de drap gris de fer galonné d’or, avec une veste rouge galonnée d’argent. ils furent au Château de St ; \textbf{Sauveur} et à \textbf{Camprieu} où ils désarmèrent le château et les soldats et bourgeoisie de Camprieu ;

Aux \textbf{Plantier} une autre troupe de fanatiques au nombre de 140 n’ayant pas la moitié d’armés, menés par le nommé Salomon lui vont joindre et sont venus deux jours avec eux ; ensuite cette troupe se retira du côté de la montagne de Lauzero. Il y a 15 jours qu’il voulait quitter Castanet, comme il le fit aux Plantiers, il promit de le rejoindre dans deux jours ; étant parti avec les nommés Etienne, fils de Fabre dit Causses Vertes de Pracoustal, le nommé Thomas d’Aumessas, le nommé Abram ne sachant d’où il est, Pompiré (?) de Ganges, un soldat qui aurait déserté son capitaine, un autre homme et un inconnu étant au nombre de sept seraient allés, portant son fusil avec de la poudre et trois balles que Laroze lui aurait baillé, des Plantiers où ils laissèrent la troupe, passèrent à \textbf{Bonnerrier} descendirent au pont del Cros passèrent à \textbf{Taleyrac} où l’un des sept fut au village acheter un pain de 2 sols qu’ils mangèrent à \textbf{Peyrefiche}, entre la vallée et \textbf{Mandagout} ils se séparèrent, le Sr. Crégut et le nommé Causses Vertes étant allés près du mas de Prascoutal paroisse d’Aulas, habitation du Sr. Causses vertes, dans la claye appelée la Borie du Gendre où ils sont restés depuis jusqu’à hier. François Martin, fournier \textit{(boulanger)} d’Aulas, que le répondant avait connu à l’assemblée qui s’était tenue lorsqu’il aidait aux châtaignes, Alibert leur ayant porté du vin, n’ayant vu plus d’autre personne. Et hier environ 9h00 du matin, le Sr. Martin et le nommé Delenne, prédicant portant un fusil et trois pistolets et le nommé Teulon de la paroisse de Valleraugue faisant le prédicant armé d’un fusil et de deux pistolets, chacun suivi de deux hommes inconnus du Sr. Crégut.

Le Sr. Teulon paraissait âgé d’environ 30 ans, gravé de la petite vérole, le visage plein, les cheveux coupés noirs, habillé de minime, de grande taille, ayant dit à Crégut et Causses Vertes qu’il fallait qu’ils vinssent avec eux pour aider à tenir une assemblée qu’il voulait faire à \textbf{Salagosse}, ce qu’ils auraient offerts.

Delenne et Teulon parlant des affaires du temps, le dit Teulon dit à Delenne qu’il était cette semaine dernière à la Coste…. Avec six de la troupe et qu’ayant vu venir les dragons et les autres troupes de la garnison du Vigan, ils auraient tiré deux coups de pistolet et ayant monté la Côte de la \textbf{Lusette}, ayant trouvé un muletier il lui aurait dit de dire aux dragons qu’ils étaient au nombre de 200, qu’ils attendaient tous les dragons montant à la Lusette. Le Sr. Teulon vit qu’ils se jetèrent dans un bois et que les dragons les abandonnèrent, et ensuite Delenne ayant posté le répondant Crégut en sentinelle près de la montagne de Salagosse où il resta de 11h du matin jusqu’au soir.

Le Sr. Martin, fournier, lui ayant porté du vin, il se tint une assemblée dans le fond de Salagosse au nombre de 200 personnes, hommes et femmes, du menu peuple de la paroisse d’Aulas. Le Sr. Delenne prêchait et plusieurs chantaient des Psaumes, le dit Crégut n’en ayant connu aucun à cause qu’il était loin d’eux.
Et le soir le Sr. Delenne serait venu relever les sentinelles. Le Sr. Delenne, le Sr. Crégut et Martin étant allés en compagnie, jusque tout près de Bréau où ayant trouvé un homme de


